\begin{Exercise}[label={ex:libretto}]
		Scrivere un programma che permetta all'utente di gestire un libretto di voti di esami, per un massimo di 20 esami. Ogni esame sostenuto ha i seguenti attributi:

		\begin{itemize}
		    \item nome del corso (stringa di massimo 256 caratteri)
		    \item voto (minimo 18, massimo 30)
		    \item data (in formato gg/mm/aaaa)
		    \item codice del corso (codice di massimo 6 cifre)
		\end{itemize}

		Il programma permette all'utente di svolgere le seguenti operazioni:

		\begin{itemize}
		    \item \textbf{inserimento:} inserisce un esame in fondo alla lista degli esami eventualmente presenti.
		    \item \textbf{stampa:} stampa tutti gli esami presenti, con i dettagli di cui sopra.
		    \item \textbf{ricerca:} chiede all'utente di inserire un codice e cerca nel libretto la presenza di un esame corrispondente a quel codice. Se presente, stampa tale esame.
		    \item \textbf{uscita:} esce dal programma.
		\end{itemize}
		tramite un menù di scelta di questo tipo:
		\begin{verbatim}
		    [i] inserimento nuovo esame
		    [s] stampa del libretto
		    [r] ricerca per codice
		    [x] uscita
		\end{verbatim}

		Il programma deve continuare a presentare tale menù, fino a quando l'utente non preme il tasto per uscire.
\end{Exercise}

\begin{Answer}[ref={ex:libretto}]
  	\lstinputlisting[style=cstylesmall]{src/libretto.c}
\end{Answer}